% Vietnamese translation for OI Public Library
% Source-PDF: 英文资料 ENGLISH MATERIALS/Longest Path on DAG.pdf
\documentclass[11pt]{article}
\usepackage{oipl-vn}

\title{Đường đi ngắn nhất nguồn đơn trong đồ thị có hướng không chu trình}
\author{Lecture Notes \#17}
\date{}

\begin{document}
\maketitle

\origtitle{Single-Source Shortest Paths in Directed Acyclic Graphs}
\sourcepath{English Materials / Longest Path on DAG.pdf}

\section{Đồ thị DAG và thứ tự tô-pô}

Trong một đồ thị có hướng không chu trình, gọi tắt là DAG, không tồn tại chu
trình. Vì vậy, chu trình âm không thể tồn tại trong DAG, và bài toán đường đi
ngắn nhất được định nghĩa tốt.

Bài toán đường đi ngắn nhất nguồn đơn trên DAG có thể giải hiệu quả hơn nhờ
sắp xếp tô-pô.

Một thứ tự tô-pô của DAG $G=(V,E)$ là một thứ tự tuyến tính của toàn bộ các
đỉnh sao cho nếu $G$ có cạnh $(u,v)$ thì $u$ xuất hiện trước $v$ trong thứ
tự. Có thể nhìn thứ tự tô-pô như cách đặt các đỉnh trên một đường ngang sao
cho mọi cạnh có hướng đều đi từ trái sang phải.

\section{Thuật toán sắp xếp tô-pô}

Sắp xếp tô-pô là một ứng dụng của tìm kiếm theo chiều sâu.

\paragraph{Thủ tục \textsc{Topsort}$(G)$.}
\begin{enumerate}
  \item Đánh dấu mọi đỉnh $v\in G$ là \emph{chưa thăm}.
  \item Khi còn tồn tại một đỉnh $v$ chưa thăm, gọi
  \textsc{Sort}$(G,v)$.
\end{enumerate}

\paragraph{Thủ tục \textsc{Sort}$(G,v)$.}
\begin{enumerate}
  \item Đánh dấu $v$ là \emph{đã thăm}.
  \item In $v$.
  \item Với mỗi $u\in \operatorname{Adj}[v]$, nếu $u$ chưa thăm thì gọi
  \textsc{Sort}$(G,u)$.
\end{enumerate}

Một cách minh họa trực quan là đặt cùng một DAG theo thứ tự tô-pô, sao cho mọi
cạnh đều đi từ trái sang phải. Độ phức tạp thời gian của \textsc{Topsort} là
\[
  O(|V|+|E|).
\]

\section{Đường đi ngắn nhất trên DAG}

Thuật toán sau giải bài toán đường đi ngắn nhất nguồn đơn trên DAG có trọng
số.

\paragraph{Thủ tục \textsc{DAG-Shortest-Paths}$(G,s,w)$.}
\begin{enumerate}
  \item Gọi \textsc{Topsort}$(G)$.
  \item Gọi \textsc{Initialize-Single-Source}$(G,s)$.
  \item Đặt $S:=\varnothing$.
  \item Xét từng đỉnh $u$ theo thứ tự tô-pô. Với mỗi $v\in\operatorname{Adj}[u]$,
  thực hiện \textsc{Relax}$(u,v,w)$, rồi đưa $u$ vào $S$.
\end{enumerate}

Hai thủ tục con giống như trong thuật toán Bellman--Ford.

\paragraph{Thủ tục \textsc{Initialize-Single-Source}$(G,s)$.}
\begin{enumerate}
  \item Với mỗi đỉnh $v\in V$, đặt $d[v]:=\infty$ và $\pi[v]:=\mathrm{nil}$.
  \item Đặt $d[s]:=0$.
\end{enumerate}

\paragraph{Thủ tục \textsc{Relax}$(u,v,w)$.}
Nếu
\[
  d[v]>d[u]+w(u,v),
\]
thì cập nhật
\[
  d[v]:=d[u]+w(u,v),\qquad \pi[v]:=u.
\]

Ví dụ trong bản gốc chạy thuật toán trên đồ thị có các đỉnh
$r,s,t,x,y,z$. Sau khi sắp theo thứ tự tô-pô, thuật toán thư giãn các cạnh từ
trái sang phải. Các giá trị cuối cùng trong ví dụ là độ dài đường đi ngắn
nhất từ nguồn $s$.

\section{Tính đúng đắn}

Định nghĩa $\delta(s,v)$ như thường lệ:
\[
  \delta(s,v)=
  \begin{cases}
    \min\{w(p): s\overset{p}{\to} v\},
      & \text{nếu tồn tại đường đi từ }s\text{ đến }v,\\
    \infty, & \text{ngược lại.}
  \end{cases}
\]

Cho một đường đi
\[
  P=v_1\to v_2\to v_3\to\cdots\to v_m.
\]
Các đỉnh trong $\{v_1,v_2,\ldots,v_{j-1}\}$ được gọi là các tiền nhiệm của
$v_j$ trên $P$.

Gọi một đường đi ngắn nhất $P$ từ $s$ đến $v$ sao cho mọi tiền nhiệm của $v$
đều nằm trong $S$ là \emph{đường đi ngắn nhất từ $s$ đến $v$ tương ứng với
$S$}.

\paragraph{Bổ đề 1.}
Gọi
\[
  (u_1,u_2,\ldots,u_n)
\]
là danh sách đỉnh theo thứ tự tô-pô, với $n=|V|$ và $u_i=s$. Ngay sau vòng
lặp ngoài thứ $j$ của \textsc{DAG-Shortest-Paths}, với mọi $k>j$, giá trị
$d[u_k]$ là trọng số đường đi ngắn nhất từ $s$ đến $u_k$ tương ứng với $S$.
Hơn nữa,
\[
  d[u_{j+1}]=\delta(s,u_{j+1}).
\]

\paragraph{Chứng minh.}
\emph{Cơ sở:} $j=i$. Sau vòng lặp thứ $i$, $S=\{s\}$ và
$d[u_k]=w(s,u_k)$; bổ đề đúng.

\emph{Giả thiết quy nạp:} giả sử bổ đề đúng với $j=m<n-1$.

\emph{Bước quy nạp:} xét $j=m+1$. Ở vòng lặp thứ $m+1$, đỉnh $u_{m+1}$ được
đưa vào $S$, và \textsc{Relax}$(u_{m+1},v,w)$ được gọi cho mọi
$v\in\operatorname{Adj}[u_{m+1}]$. Do thứ tự tô-pô, các đỉnh này đều nằm bên
phải $u_{m+1}$.

Nếu $u_{m+1}$ không đạt được từ $s$, thì mọi $v$ được xét ở đây cũng không
đạt được từ $s$ qua $u_{m+1}$, và giá trị $d[v]$ vẫn là $\infty$ khi không có
đường khác.

Nếu $u_{m+1}$ đạt được từ $s$, thì với cạnh $(u_{m+1},v)$, phép thư giãn kiểm
tra liệu đường đi qua $u_{m+1}$ có tốt hơn đường đi ngắn nhất hiện tại tương
ứng với tập tiền nhiệm cũ hay không. Khi
\[
  d[v]>d[u_{m+1}]+w(u_{m+1},v),
\]
ta cập nhật
\[
  d[v]:=d[u_{m+1}]+w(u_{m+1},v).
\]
Theo giả thiết quy nạp, giá trị mới này là trọng số của đường đi ngắn nhất từ
$s$ đến $v$ tương ứng với tập mới
\[
  S=\{u_i,u_{i+1},\ldots,u_{m+1}\}.
\]

Ngoài ra, nếu $v=u_{m+2}$ thì không có đỉnh nào trong
\[
  \{u_{m+3},u_{m+4},\ldots,u_n\}
\]
có thể đi đến $u_{m+2}$, vì thứ tự tô-pô đặt mọi cạnh từ trái sang phải. Do
đó
\[
  d[v]=\delta(s,v).
\]
Bổ đề đúng cho vòng lặp thứ $m+1$, hoàn tất quy nạp.

\paragraph{Định lý 1.}
Thuật toán \textsc{DAG-Shortest-Paths}, chạy trên DAG có trọng số
$G=(V,E)$ với nguồn $s$, kết thúc với
\[
  d[v]=\delta(s,v)
\]
cho mọi đỉnh $v\in V$.

\paragraph{Chứng minh.}
Theo Bổ đề 1, sau vòng lặp thứ $j$, giá trị
\[
  d[u_{j+1}]=\delta(s,u_{j+1})
\]
đã được chốt. Sau $n-1$ vòng lặp, mọi đỉnh đều có giá trị đúng. Vòng lặp
ngoài thật ra chạy $n$ lần; vòng cuối cùng là dư thừa nhưng không làm thay
đổi tính đúng đắn.

\section{Phân tích độ phức tạp}

\begin{itemize}
  \item \textsc{Topsort} tốn $O(|V|+|E|)$ thời gian.
  \item \textsc{Initialize-Single-Source} tốn $O(|V|)$ thời gian.
  \item Vòng lặp ngoài chạy $|V|$ lần. Tuy vậy, xét tổng toàn bộ vòng lặp
  trong, mỗi cạnh chỉ được duyệt đúng một lần từ trái sang phải, nên tổng số
  lần lặp bên trong là $|E|$. Mỗi lần lặp bên trong tốn $O(1)$.
\end{itemize}

Vì vậy tổng thời gian chạy là
\[
  O(|V|+|E|).
\]

\section{Ghi chú về đường đi dài nhất}

Với DAG, bài toán đường đi dài nhất cũng có thể xử lý bằng cùng ý tưởng thứ
tự tô-pô, vì không có chu trình để làm đường đi tăng vô hạn. Một cách nhìn
đơn giản là đổi dấu trọng số cạnh rồi chạy thuật toán đường đi ngắn nhất trên
DAG; hoặc trực tiếp thay phép thư giãn nhỏ nhất bằng phép thư giãn lớn nhất.
Thứ tự tô-pô vẫn bảo đảm khi xét một đỉnh, mọi tiền nhiệm của nó đã được xử
lý.

\end{document}
